
\ifdefined\FULLARTICLE
\else
  \documentclass[12pt,a4paper]{article}
  \usepackage[T2A]{fontenc}
  \usepackage[utf8]{inputenc}
  \usepackage[russian]{babel}
  \usepackage{amsmath, amssymb, amsthm}
  \usepackage{geometry}
  \usepackage{booktabs} % для красивых таблиц
  \usepackage{tabularx}
  \geometry{margin=2.5cm}

  % Настройка окружений теорем (независимые счётчики для корректной нумерации)
  \theoremstyle{definition}
  \newtheorem{definition}{Definition}
  \newtheorem{assumption}{Assumption}
  \newtheorem{lemma}{Lemma}
  \newtheorem{theorem}{Theorem}
  \newtheorem{proposition}{Proposition}
  \newtheorem{corollary}{Corollary}
  \newtheorem{observation}{Observation}
  \newtheorem{remark}{Remark}
  \renewcommand{\proofname}{Proof}

  \begin{document}

  \begin{center}
    \Large\textbf{Phase-Based Planning Model for Software Work in a Human--Agent Cell}
  \end{center}
\fi

\subsection{Formal Definitions and Deterministic Assumptions}

\begin{definition}\label{def:params}
Пусть заданы следующие параметры:
\begin{itemize}
  \item $X > 0$ --- базовая трудоёмкость единицы работы (время реализации элементарной функциональности);
  \item $Z_i \in \mathbb{R}^{+}$ --- сложность $i$-й задачи, измеренная в единицах $X$. Дробные значения допустимы: при детализации проекта подзадача может занимать долю базовой единицы работы (например, $Z_i = 0.5$). В прикладном расчёте значения обычно выбираются из конечной дискретной шкалы $D\subset\mathbb{Q}^{+}$;
  \item $r \in R$ --- \textbf{режим} организации работы с ИИ, то есть способ взаимодействия разработчика с инструментом (например, интерактивная работа в одном диалоге, делегированное выполнение с автономной агентной фазой и последующей проверкой либо полностью автономное выполнение без обязательной human-фазы). Основной анализ M4 относится к интерактивному и делегированному режимам; полностью автономный режим служит предельным случаем. Конечное множество режимов $R$ задаётся при применении модели;
  \item $k_i^{(r)} \in \mathbb{R}^+$ --- нормированная длительность $i$-й задачи при использовании ИИ в режиме $r$. Величина является свойством тройки «задача $\times$ инструмент $\times$ режим» и определяется операционально:
  \[
    k_i^{(r)} = \frac{T_i^{ai}(r)}{Z_i X},
  \]
  где $T_i^{ai}(r)$ --- фактическое время выполнения $i$-й задачи с применением ИИ \emph{в одном рабочем потоке режима $r$ при полной доступности разработчика}, то есть без конкуренции с другими потоками за его внимание. Оно включает формулирование запросов, ожидание генерации, проверку и исправление результата. Значение $k_i^{(r)} < 1$ соответствует более короткой длительности относительно базовой оценки $Z_iX$, $k_i^{(r)} = 1$ --- совпадению с ней, а $k_i^{(r)} > 1$ --- более длинной длительности. Это описательная нормировка, а не причинная оценка эффекта ИИ без отдельного идентификационного дизайна. Если режим зафиксирован или ясен из контекста, верхний индекс опускается: $k_i := k_i^{(r)}$;
  \item $h_i^{(r)},\, a_i^{(r)}$ --- \textbf{человеческая} и \textbf{агентная} составляющие коэффициента $k_i^{(r)}$. Для определения $k_i^{(r)}$ время $T_i^{ai}(r)$ измеряется по двум каналам. К первому относятся интервалы, в течение которых разработчик занят задачей: формулирует запросы, читает и проверяет результаты, направляет исправления. Ко второму относятся интервалы автономной работы агента, когда разработчик свободен: генерация, самопроверка и выполнение тестов. Внутри задачи эти интервалы могут чередоваться; их порядок задаётся фазами уровня M4. После нормирования обеих частей на $Z_i X$ по построению получаем:
  \[
    k_i^{(r)} = a_i^{(r)} + h_i^{(r)}, \qquad 0 \le h_i^{(r)} \le k_i^{(r)}.
  \]
  Текущая модель относит каждый учитываемый интервал ровно к одному из этих двух каналов; отдельные ограниченные ресурсы CI, API и внешней инфраструктуры в неё не входят. При $P = 1$ это разбиение не влияет на длительность задачи: без очереди к человеку она равна $k_i^{(r)} Z_i X$ при любом соотношении составляющих. Различие становится существенным при $P > 1$: агентные этапы разных задач могут выполняться параллельно, тогда как все человеческие этапы обслуживает один разработчик (уровень M4, Замечание~\ref{rem:human});
  \item $P \ge 1$ --- настроенное число одновременно доступных агентных потоков; для исходного сценария без ИИ принимается $P_h = 1$. На уровнях M1--M4, где задачи и фазы неделимы, $P$ является целым. Только в агрегированных формулах M0 величину $P$ допустимо интерпретировать как эффективный, усреднённый за период параллелизм. Важно различать дополнительную \emph{доступную} мощность и реально выбранную конфигурацию: неиспользуемый поток сам по себе не замедляет расписание, тогда как переход к конфигурации с большим $P$ может изменить $C(P)$ и $\gamma(P)$;
  \item $\sigma$ --- \textbf{конфигурация процесса}: пара $\sigma = (P, r)$ либо, в общем случае, пара $(P, \rho)$, где $\rho\colon \{1, \dots, N\} \to R$ --- назначение режимов задачам (разные потоки работы могут вестись в разных режимах). Эта запись является сокращённой при фиксированных ИИ-инструменте и правилах приёмки: если инструмент или критерий принятого результата меняются, их необходимо явно указывать при сравнении, хотя они не добавляются в размерность $\sigma$. Производные величины модели --- суммарная работа $W$, средневзвешенный коэффициент $\bar{k}_w$, время $T_{ai}$ и ускорение $S_E$ --- являются функциями конфигурации. Сравнение «вариантов организации работы» в модели означает сравнение конфигураций целиком: параметры $P$ и $\{k_i^{(r)}\}$ не варьируются независимо друг от друга (см. Замечание~\ref{rem:factorization});
  \item $N \in \mathbb{N}$, $N\ge1$ --- общее число задач в проекте;
  \item $T_h,\; T_{ai}$ --- общее время выполнения всех задач без ИИ и с ИИ соответственно.
\end{itemize}
\end{definition}

\begin{remark}[Task Independence]\label{rem:independent}
Сначала предполагается, что все задачи $i = \overline{1,N}$ \textbf{не пересекаются}, то есть не используют общие ресурсы, не блокируют друг друга и могут выполняться параллельно без дополнительных затрат на синхронизацию. Это допущение позволяет использовать аддитивную модель времени. Зависимости между задачами и критический путь вводятся далее на уровне M2.
\end{remark}

\subsection{Model Hierarchy Overview}

\begin{remark}[Model Hierarchy M0--M4]\label{rem:hierarchy}
Уровни M0--M4 образуют последовательность уточнений одной постановки. Каждый
следующий уровень снимает часть релаксаций и вводит одну или несколько
связанных сторон модели; это не означает, что переход между любыми двумя
соседними уровнями добавляет ровно одно ограничение.

\begin{center}
\small
\begin{tabularx}{\textwidth}{@{}c>{\raggedright\arraybackslash}X>{\raggedright\arraybackslash}p{0.42\textwidth}@{}}
\toprule
Уровень & Основное уточнение & Нижняя оценка \\
\midrule
M0 & идеально делимая агрегированная работа & $B_0 = W/P$ \\
M1 & неделимость задач & $B_1 = \max\{W/P,\; M_N\}$ \\
M2 & граф зависимостей & $B_2 = \max\{W/P,\; L\}$ \\
M3 & эффективные веса и интеграционные издержки & $B_3 = \max\{W_3/P,\; L_3\}$ \\
M4 & фазы, локальное закрепление, очередь и общий человек & $B_4 = \max\{W_4/P,\; L_4,\; \gamma H\}$ \\
\bottomrule
\end{tabularx}
\end{center}

Формальные уточнения вводятся соответственно в
Предположении~\ref{assum:ideal} и
Замечаниях~\ref{rem:indivisible}--\ref{rem:human}.
\end{remark}

\subsection{Level M0: Ideally Divisible Work}

\begin{lemma}[Baseline Time without AI]\label{lem:baseline}
При последовательном выполнении задач разработчиком без использования ИИ общее время составляет:
\begin{equation}
  T_h = \sum_{i=1}^{N} Z_i X.
\end{equation}
\end{lemma}

\begin{assumption}[Idealized AI-Assisted Time Model (Level M0)]\label{assum:ideal}
В рамках основной аналитической модели работа считается идеально делимой, а параллелизм распределён равномерно. Обозначим суммарный объём работы с ИИ:
\[
  W := X \sum_{i=1}^{N} Z_i k_i.
\]
При использовании ИИ с коэффициентом параллелизма $P$ общее время выполнения задач принимается равным:
\begin{equation}
  T_{ai}^{(0)} = \frac{W}{P} = \frac{1}{P} \sum_{i=1}^{N} Z_i k_i X.
\end{equation}
Все дальнейшие теоретические выводы строятся на этом допущении, если явно не указан иной уровень иерархии моделей (Замечание~\ref{rem:hierarchy}).
\end{assumption}

\begin{observation}[Speedup Criterion and Factor at Level M0]\label{obs:ideal}
Обозначим средневзвешенную нормированную длительность (веса --- сложности задач $Z_i$):
\[
  \bar{k}_w := \frac{\sum_{i=1}^{N} Z_i k_i}{\sum_{i=1}^{N} Z_i}.
\]
Подстановка Леммы~\ref{lem:baseline} и Предположения~\ref{assum:ideal} непосредственно даёт цепочку эквивалентностей
\begin{equation}
  T_{ai}^{(0)} < T_h
  \quad\Longleftrightarrow\quad
  \sum_{i=1}^{N} Z_i k_i < P \sum_{i=1}^{N} Z_i
  \quad\Longleftrightarrow\quad
  P > \bar{k}_w.
\end{equation}
Это не самостоятельный теоретический результат, а алгебраическая перезапись определения времени на уровне M0.

Коэффициент ускорения равен
\begin{equation}
  S_E = \frac{T_h}{T_{ai}^{(0)}} = P \cdot \frac{\sum_{i=1}^{N} Z_i}{\sum_{i=1}^{N} Z_i k_i} = \frac{P}{\bar{k}_w}.
\end{equation}
В частности, если $k_i \le 1$ для всех $i$ и $P > 1$, то $S_E \ge P > 1$, т.\,е. ускорение гарантировано.

\medskip
Кроме того, из $\min_{j} k_j \le \bar{k}_w \le \max_{j} k_j$ следует двусторонняя оценка
\[
  \frac{P}{\max_{i} k_i} \;\le\; S_E \;\le\; \frac{P}{\min_{i} k_i},
\]
позволяющая быстро оценивать диапазон ускорения без взвешивания по $Z_i$.
\end{observation}

\begin{remark}[Speedup Factorization and Baseline Choice]\label{rem:factorization}
Формула Наблюдения~\ref{obs:ideal} допускает факторизацию
\[
  S_E = \underbrace{P\vphantom{\frac{1}{\bar{k}_w}}}_{\text{организационный множитель}} \cdot \underbrace{\frac{1}{\bar{k}_w}}_{\text{инструментальный множитель}}.
\]
Первый множитель отражает эффект параллельной организации работы, а второй --- влияние ИИ на трудоёмкость задач. В качестве исходного сценария $T_h$ намеренно выбрана последовательная работа одного разработчика ($P_h = 1$). Тем самым модель отвечает на вопрос о том, что получает один разработчик при переходе к ИИ-агентам, а не сравнивает ИИ с командой из $P$ человек. При $k_i \equiv 1$ получаем $S_E = P$: такое ускорение целиком обусловлено организацией процесса и в принципе достижимо также командой из $P$ людей. Основное допущение состоит в том, что $P$ параллельных потоков обеспечивают ИИ-агенты под контролем одного и того же разработчика, без привлечения дополнительных специалистов. Издержки контроля учитываются на уровнях M3--M4: функция $C(P)$ описывает межагентную интеграцию (Замечание~\ref{rem:coordination}), а член $\gamma(P)\,H$ --- человеческий ресурс и связанный с ним потолок $S_E^{\mathrm{ach}} \le 1/\bar{h}_w$ (Замечание~\ref{rem:human}, Предложение~\ref{prop:amdahl}). Для сравнения с командой из $q$ разработчиков исходное время следует заменить на $T_h(q) = T_h/q$ в идеализированном приближении; тогда условие превосходства процесса с ИИ принимает вид $P/\bar{k}_w > q$.

\medskip
Факторизация справедлива \emph{для каждой фиксированной конфигурации} $\sigma = (P, r)$, но её множители нельзя варьировать независимо. На практике изменение $P$ часто сопровождается сменой режима $r$, например переходом от интерактивной работы к автономным агентам. Вместе с режимом меняется весь вектор $\{k_i^{(r)}\}$ и, следовательно, $\bar{k}_w$. Поэтому нельзя оценить $\bar{k}_w$ в одном режиме и использовать параллелизм другого, например измерить $k_i$ при интерактивной работе и подставить $P = 4$ для автономных агентов. Сравнивать следует конфигурации целиком: $S_E(\sigma_1)$ и $S_E(\sigma_2)$.
\end{remark}

\begin{remark}[Critical Parallelism Threshold for $k_i > 1$]\label{rem:threshold}
Из тождества Наблюдения~\ref{obs:ideal}
\[
  \sum_{i=1}^{N} Z_i k_i < P \sum_{i=1}^{N} Z_i
\]
следует, что ускорение $S_E > 1$ достигается тогда и только тогда, когда
\[
  P > \bar{k}_w,
\]
где $\bar{k}_w$ --- средневзвешенная нормированная длительность, введённая в Наблюдении~\ref{obs:ideal}.

Подчеркнём: условие $P > \bar{k}_w$ необходимо и достаточно только на уровне M0. На уровнях M1--M2 оно остаётся необходимым, но перестаёт быть достаточным: достижимое время дополнительно ограничено снизу самой длинной задачей ($M_N$, Замечание~\ref{rem:indivisible}) и критическим путём ($L$, Замечание~\ref{rem:graph}); содержательный достаточный результат для этих уровней даёт Теорема~\ref{thm:sufficient}. На уровне M3 необходимое условие по суммарной занятости потоков принимает вид $C(P)\bar a_w+\bar h_w<P$ (Замечание~\ref{rem:coordination}). На уровне M4 необходимы одновременно все три неравенства, задаваемые $B_4<T_h$: $W_4/P<T_h$, $L_4<T_h$ и $\gamma(P)H<T_h$. Они не являются достаточными, поскольку существуют экземпляры с $T^{*}>B_4$ из-за локальных блокировок и очереди к человеку.

\begin{itemize}
  \item Если $k_i \equiv k > 1$ для всех $i$ (ИИ равномерно замедляет выполнение), то критический порог упрощается до:
  \[
    P_{\text{crit}} = k.
  \]
  То есть параллелизм должен \textbf{строго превышать} коэффициент замедления, иначе использование ИИ нецелесообразно.

  \item Если коэффициенты $k_i$ неоднородны, то порог определяется взвешенным средним: задачи большей сложности вносят больший вклад в $\bar{k}_w$. Это означает, что даже при наличии отдельных задач с $k_i \gg 1$ ускорение возможно, если их вклад в общую сложность мал, а параллелизм достаточен для компенсации.

  \item В практическом применении $\bar{k}_w$ характеризует длительность выбранной AI-поддержанной конфигурации для данного проекта. Процедура оценки приведена в Замечании~\ref{rem:estimation}. Если $\bar{k}_w$ не меньше доступного $P$, эта конфигурация не даёт временного выигрыша. В этом случае следует:
  \begin{enumerate}
    \item пересмотреть состав автоматизируемых задач, начиная с задач с наибольшими $k_i$;
    \item увеличить допустимый параллелизм процесса, например за счёт асинхронного выполнения;
    \item изменить способ применения ИИ --- постановку запросов, состав контекста и последующую обработку --- так, чтобы снизить $k_i$.
  \end{enumerate}
\end{itemize}

Итак, модель задаёт не бинарную оценку полезности ИИ, а \textbf{количественное условие принятия решения}, которое связывает длительности в выбранном режиме ($k_i^{(r)}$), структуру проекта ($Z_i$) и доступный команде параллелизм ($P$). Критерий применяется к конфигурации $\sigma = (P, r)$ целиком: значения $\bar{k}_w$ и $P$ должны относиться к одному режиму работы (Замечание~\ref{rem:factorization}). Выбор между полностью ручным и AI-поддержанным выполнением отдельных задач лежит вне постановки; после такого выбора модель планирует только включённый AI-поддержанный набор.
\end{remark}

\subsection{Level M1: Task Indivisibility}

\begin{remark}[Workload Lower Bound $B_1$]\label{rem:indivisible}
Формула уровня M0 предполагает идеальное равномерное распределение нагрузки. В реальности для неделимых задач время выполнения не может быть меньше длительности самой длинной задачи, даже при сколь угодно большом $P$. Обозначим
\[
  M_N := X \max_{i \in \overline{1,N}} (Z_i k_i).
\]
Более строгая оценка записывается как неравенство:
\[
  T_{ai} \ge B_1 := \max\left\{ \frac{W}{P},\; M_N \right\}.
\]
Здесь $B_1$ --- нижняя оценка времени завершения всех работ; точное значение зависит от распределения задач между целым числом исполнителей и может строго превышать $B_1$ даже при оптимальном расписании.
\end{remark}

\subsection{Level M2: Dependencies and the Critical Path}

\begin{remark}[Critical-Path Lower Bound $B_2$]\label{rem:graph}
При наличии логических или ресурсных зависимостей вместо аддитивной модели используется графовая. Пусть $G = (V, E)$ --- ориентированный ациклический граф, вершины $V$ которого соответствуют задачам, а дуги $E$ задают отношения предшествования. Тогда общее время выполнения с ИИ удовлетворяет неравенству:
\[
  T_{ai} \ge B_2 := \max\left\{ \frac{W}{P},\; L \right\},
\]
где $L$ --- длина критического пути в $G$ с весами $Z_i k_i X$. Любая отдельная задача образует путь в $G$, поэтому $L \ge M_N$ и уровень M2 обобщает M1: при пустом множестве дуг ($E = \varnothing$) выполняется $L = M_N$ и $B_2 = B_1$. Равенство $T_{ai} = B_2$ в общем случае не гарантируется, так как задача планирования на ограниченном числе исполнителей является NP-трудной.
\end{remark}

\subsubsection{Attainability and Speedup Guarantees for M1--M2}

\begin{remark}[Attainability Guarantee: Graham's Bound]\label{rem:graham}
На уровнях M1--M2, при целом числе $P$ одинаковых исполнителей и $C(P) \equiv 1$, нижние оценки $B_1$ и $B_2$ можно дополнить гарантией достижимости. Из классического результата Грэма для списочного планирования с фиксированным приоритетом, при котором исполнитель не простаивает, если имеется готовая задача, следует неравенство \cite{graham1966bounds}
\[
  T_{ai} \;\le\; \frac{W}{P} + \left(1 - \frac{1}{P}\right) L \;\le\; \left(2 - \frac{1}{P}\right) B_2.
\]
Для независимых задач (уровень M1) те же соотношения выполняются с заменой $L$ на $M_N$ и $B_2$ на $B_1$. Таким образом, оптимальное время завершения $T^{*}$ находится в интервале
\[
  B_j \;\le\; T^{*} \;\le\; \left(2 - \frac{1}{P}\right) B_j, \qquad j \in \{1, 2\}.
\]
Следовательно, достижимое ускорение отличается от верхней оценки $\bar S^{(j)} = T_h/B_j$ не более чем в $2 - 1/P$ раза:
\[
  \frac{\bar S^{(j)}}{2 - 1/P} \;\le\; S_E^{\text{ach}} \;\le\; \bar S^{(j)}.
\]
Для независимых задач при упорядочивании по невозрастанию длительностей (правило LPT) известна более сильная гарантия $T_{\mathrm{LPT}} \le \left(4/3 - 1/(3P)\right) T^{*}$ \cite{graham1969bounds}. Однако эта константа относится к оптимуму $T^{*}$, а не к нижней границе $B_j$, и потому не переносится на приведённый выше интервал. Отношение $T^{*}/B_1$ может превышать $4/3 - 1/(3P)$: например, для трёх задач равной длительности при $P = 2$ имеем $T^{*}/B_1 = 4/3 > 4/3 - 1/6$. На уровни M3--M4 эти гарантии непосредственно не распространяются, поскольку постановка Грэма не включает координационные издержки и общий человеческий ресурс. Такие ресурсы рассматриваются в более общих моделях теории расписаний, в частности в RCPSP и задачах параллельных машин с одним общим обслуживающим устройством \cite{brucker1999resource,hall2000parallel,kravchenko1997parallel}.
\end{remark}

\begin{theorem}[Sufficient Speedup Criterion at Levels M1--M2]\label{thm:sufficient}
Пусть $C(P) \equiv 1$, $P$ --- целое число одинаковых исполнителей. Обозначим относительную длину узкого места:
\[
  \mu := \frac{M_N}{T_h} \quad \text{(уровень M1)}, \qquad \mu := \frac{L}{T_h} \quad \text{(уровень M2)}.
\]
Если
\begin{equation}\label{eq:sufficient}
  \bar{k}_w < P - (P - 1)\, \mu,
\end{equation}
то ускорение гарантировано: для любого жадного расписания выполняется $T_{ai} < T_h$. В более общем случае целевой уровень ускорения $S_E^{\text{ach}} \ge S_{\text{target}} \ge 1$ гарантирован при
\begin{equation}\label{eq:sufficient-target}
  \bar{k}_w \;\le\; \frac{P}{S_{\text{target}}} - (P - 1)\, \mu.
\end{equation}
При $S_{\text{target}} = 1$ нестрогая версия условия~\eqref{eq:sufficient-target}
гарантирует отсутствие замедления, $S_E^{\mathrm{ach}}\ge1$, тогда как её
строгая версия совпадает с~\eqref{eq:sufficient} и гарантирует собственно
ускорение, $S_E^{\mathrm{ach}}>1$. При $\mu \to 0$ (отсутствие доминирующего
узкого места) строгий достаточный критерий сходится к необходимому
$\bar{k}_w < P$ из Замечания~\ref{rem:threshold}. Зазор между необходимым и
достаточным условиями равен $(P-1)\,\mu$ --- это «плата за неделимость» (M1)
или «плата за зависимости» (M2).
\end{theorem}

\begin{proof}
По Замечанию~\ref{rem:graham} любое жадное расписание удовлетворяет $T_{ai} \le W/P + (1 - 1/P)\,L$ (на уровне M1 --- с $M_N$ вместо $L$). Разделив обе части на $T_h$ и учитывая $W = \bar{k}_w\,T_h$, получаем
\[
  \frac{T_{ai}}{T_h} \;\le\; \frac{\bar{k}_w}{P} + \left(1 - \frac{1}{P}\right) \mu.
\]
Правая часть не превосходит $1/S_{\text{target}}$ в точности при условии~\eqref{eq:sufficient-target} (умножение на $P$ и перенос слагаемого), откуда $S_E^{\text{ach}} = T_h/T_{ai} \ge S_{\text{target}}$; случай~\eqref{eq:sufficient} получается при $S_{\text{target}} = 1$ со строгим неравенством. Ч.Т.Д.
\end{proof}

\begin{remark}[Equivalent Form via Bottleneck Dominance]\label{rem:dominance}
Введём относительную доминантность узкого места --- отношение самой длинной задачи (на M2 --- критического пути) к средней нагрузке на исполнителя:
\[
  b := \frac{M_N}{W/P} \quad \text{(M1)}, \qquad b := \frac{L}{W/P} \quad \text{(M2)}.
\]
Параметры связаны точным соотношением $\mu = b\,\bar{k}_w/P$, поэтому критерий~\eqref{eq:sufficient} эквивалентно записывается в виде
\[
  \bar{k}_w \left(1 + \left(1 - \frac{1}{P}\right)b\right) < P.
\]
Одно неравенство выполняется тогда и только тогда, когда выполняется другое, --- это тот же критерий в координатах $(\bar{k}_w, b)$ вместо $(\bar{k}_w, \mu)$.

Из формы через $b$ следуют два упрощённых практических условия. Они \emph{не эквивалентны} исходному критерию, но дают более консервативную оценку: каждое условие строго сильнее~\eqref{eq:sufficient}, поэтому его выполнение по-прежнему гарантирует ускорение.
\begin{enumerate}
  \item \textbf{Упрощение коэффициента.} Если заменить множитель $1 - 1/P < 1$ единицей, получим условие
  \[
    \bar{k}_w\,(1 + b) < P,
  \]
  поправка которого не зависит от $P$ при фиксированном $b$. Из него следует~\eqref{eq:sufficient}, но не наоборот: критерий~\eqref{eq:sufficient} может выполняться и при $\bar{k}_w(1+b) \ge P$. При сравнении конфигураций $P$ величину $b=b(P)$ следует пересчитывать совместно с $W/P$.
  \item \textbf{Фиксация $b$ как правила декомпозиции.} Если разбиение работ гарантирует $b \le b_0$, то есть ни одна задача и ни один путь зависимостей не превышают $b_0$ средних нагрузок на исполнителя, то условие
  \[
    \bar{k}_w\,(1 + b_0) < P
  \]
  достаточно для ускорения независимо от деталей конкретного набора задач. Необходимым оно не является: фактическая доминантность узкого места может оказаться меньше $b_0$, и ускорение возможно даже при нарушении условия. Отсюда следует практическая рекомендация: декомпозировать работу до $b \le b_0$ и обеспечивать запас параллелизма $P > \bar{k}_w(1 + b_0)$.
\end{enumerate}
\end{remark}

\subsection{Level M3: Coordination Overhead}

\begin{remark}[Effective Weights and Lower Bound $B_3$]\label{rem:coordination}
Для учёта межагентной интеграции и повторной работы, возникающей при параллельном изменении общих артефактов, вводится функция $C(P) \ge 1$. Предполагается, что $C(1)=1$ и $C(P)$ не убывает с ростом $P$; для каждого конкретного процесса это предположение требует эмпирической проверки. Функция применяется только к автономной агентной работе, но не к человеческим фазам. В фиксированной конфигурации $C(P)$ считается экзогенной константой: она не зависит от назначения задач, ресурсного порядка и реализовавшихся очередей. Обозначим
\[
  A := X\sum_{i=1}^{N}Z_i a_i, \qquad
  H := X\sum_{i=1}^{N}Z_i h_i, \qquad W=A+H.
\]
При смешанном назначении режимов здесь и далее под $a_i$ и $h_i$ понимаются
$a_i^{(\rho(i))}$ и $h_i^{(\rho(i))}$.
Эффективная длительность задачи и суммарная занятость агентных потоков на уровне M3 равны
\[
  w_i^{(3)}(P):=Z_iX\bigl(C(P)a_i+h_i\bigr), \qquad
  W_3(P):=\sum_i w_i^{(3)}(P)=C(P)A+H.
\]
Пусть $L_3(P)$ --- длина критического пути в $G$ с весами $w_i^{(3)}(P)$. Тогда
\[
  T_{ai} \ge B_3 := \max\left\{ \frac{W_3(P)}{P},\; L_3(P) \right\}.
\]
Такое определение устраняет асимметрию прежней агрегированной записи: дополнительные интеграционные затраты увеличивают не только суммарную работу, но и те пути, на которых они фактически возникают.

Условие ускорения по линейной части в модели с координационными издержками принимает вид:
\[
  \frac{C(P)A+H}{P} < T_h
  \;\iff\;
  C(P)\,\bar a_w+\bar h_w<P,
\]
где $\bar a_w:=\sum_iZ_i a_i/\sum_iZ_i$ и $\bar h_w:=\sum_iZ_i h_i/\sum_iZ_i$.

Чтобы разграничить $k_i^{(r)}$ и издержки параллелизма, примем следующее правило учёта. В $k_i^{(r)}$ входит всё, что можно измерить при $P = 1$ в режиме $r$: формулирование запросов, ожидание генерации, проверка и исправление результата в одном потоке. К издержкам параллелизма относятся только эффекты, возникающие при $P > 1$. Без этого правила одни и те же затраты могут быть учтены дважды --- и в $k$, и в множителях при $P > 1$ --- либо, напротив, остаться неучтёнными, если координационные издержки неявно включены в измеренное $k$, а дополнительные множители приняты равными единице.

С учётом разложения $k_i^{(r)} = a_i^{(r)} + h_i^{(r)}$ (Определение~\ref{def:params}) и звездообразной топологии «один разработчик --- $P$ агентов» издержки разделяются на три группы:
\begin{enumerate}
  \item координационные издержки \emph{на уровне отдельной задачи} --- подготовка спецификации, передача контекста и проверка результата --- не зависят от $P$ и уже входят в $k_i^{(r)}$; поэтому они увеличивают веса задач, в том числе на критическом пути;
  \item издержки \emph{разработчика}, растущие с $P$, --- переключение между потоками и восстановление контекста при обслуживании запросов --- относятся к человеческому ресурсу и учитываются множителем $\gamma(P)$ при $H$ на уровне M4 (Замечание~\ref{rem:human});
  \item функция $C(P)$ описывает \emph{межагентные} интеграционные издержки, возрастающие с числом параллельных потоков: конфликты слияния, рассогласование интерфейсов и повторную работу из-за одновременных изменений.
\end{enumerate}
В качестве рабочей параметрической формы можно использовать знаменатель универсального закона масштабируемости Гантера:
$C(P) = 1 + \alpha (P - 1) + \beta P (P - 1)$, где $\alpha,\beta\ge0$, $\alpha$ характеризует конкуренцию за ресурсы, а $\beta$ --- затраты на согласование \cite{gunther2008general}. Применение этой зависимости к агентам программирования является гипотезой настоящей модели, а не эмпирическим результатом указанной работы. Формула Брукса с $\sim P(P-1)/2$ попарными каналами связи относится к человеческой команде и не описывает звездообразную топологию буквально. Однако косвенная связанность через общий код, интерфейсы и интеграцию сохраняется даже без непосредственного общения агентов \cite{brooks1995mythical}.
\end{remark}

\subsection{Level M4: Phase Semantics and the Exact Schedule}

\begin{remark}[Local Blocking and Lower Bound $B_4$]\label{rem:human}
Разложение $k_i^{(r)} = a_i^{(r)} + h_i^{(r)}$ (Определение~\ref{def:params}) на уровне M4 дополняется порядком фаз. Задача $i$ представляется конечной последовательностью $\mathcal F_i=(f_{i1},\dots,f_{im_i})$. Агентные фазы включают автономную реализацию, самопроверку и повторную работу; человеческие --- постановку, уточнение требований, проверку и приёмку. Суммы их длительностей при $P=1$ равны соответственно $Z_iXa_i$ и $Z_iXh_i$. Фазы одной задачи следуют в заданном порядке, поэтому модель допускает несколько запросов к человеку и повторные циклы «проверка --- исправление --- повторная проверка».

После начала первой фазы задача закрепляется за одним агентом до принятия результата. Человеческая фаза требует одновременно внимания разработчика и сохранения выделенного агентного потока. Если разработчик занят, запрос встаёт в очередь, а соответствующий агент прекращает работу и остаётся заблокированным. Остальные агенты продолжают свои задачи независимо, пока сами не запросят человеческое внимание. Таким образом, блокировка является \emph{локальной}, а не глобальной.

В каждый момент незавершённые назначенные задачи занимают не более $P$ агентных потоков, а разработчик обслуживает не более одной человеческой фазы. Переиспользование заблокированного потока для другой задачи и запуск дополнительного, приоритетного потока сверх $P$ в модели не допускаются. Новая задача может быть назначена только свободному агенту.
Если $(i,j)\in E$, первая фаза задачи $j$ может начаться только после приёмки
задачи $i$. Агентные и человеческие фазы считаются непрерываемыми; порядок
обслуживания одновременно ожидающих запросов является частью расписания $\pi$.

При $P>1$ эффективная длительность человеческих фаз принимается равной $\gamma(P)H$, где $\gamma(P)\ge1$, $\gamma(1)=1$. Множитель описывает дополнительное время обслуживания из-за переключения и восстановления контекста, но не включает ожидание в очереди. Как и $C(P)$, для фиксированной конфигурации он считается экзогенным и не зависит от конкретного расписания. В качестве редуцированного приближения используется $\gamma(P)=1+\delta(P-1)$, $\delta\ge0$; при эмпирической калибровке его следует связывать с фактическим числом одновременно сопровождаемых потоков и частотой вмешательств.

Определим эффективные веса и агрегаты уровня M4:
\[
  w_i^{(4)}(P):=Z_iX\bigl(C(P)a_i+\gamma(P)h_i\bigr),
\]
\[
  W_4(P):=\sum_iw_i^{(4)}(P)=C(P)A+\gamma(P)H,
\]
а через $L_4(P)$ обозначим длину критического пути с весами $w_i^{(4)}(P)$. Тогда фиксированная нижняя оценка времени имеет вид
\[
  T^{*} \;\ge\; B_4 := \max\left\{ \frac{W_4(P)}{P},\; L_4(P),\; \gamma(P)H \right\},
\]
где $T^{*}:=\min_{\pi}T(\pi)$ --- оптимальное время среди допустимых расписаний, а $T(\pi)$ --- время завершения конкретного расписания $\pi$.

Пусть $Q(\pi)$ --- суммарное время, в течение которого уже назначенные агенты заблокированы в очереди к разработчику. Это время возникает из самого расписания и не входит ни в $k_i$, ни в $H$. Для любого допустимого расписания дополнительно выполняется
\[
  T(\pi)\ge \frac{W_4(P)+Q(\pi)}{P}.
\]
Поэтому равенство $T^{*}=B_4$ не гарантируется: одновременные запросы к человеку могут увеличить срок даже при небольшом $H$.

Например, если из двух параллельно работающих агентов агент A запросил помощь, он останавливается и удерживает свой поток до обслуживания. Агент Б при этом продолжает работу. Если затем помощь потребуется и ему, оба агента окажутся в очереди к одному разработчику; простой всей системы возникнет как следствие совпадения запросов, а не как правило глобальной блокировки.

Агрегаты $A$, $H$, $W_4$, $L_4$ и $B_4$ описывают объём работы и три обязательных ограничения, но не кодируют взаимное расположение человеческих фаз. Поэтому два экземпляра с одинаковыми значениями этих агрегатов могут иметь разные $T^{*}$: нижняя граница одинакова, а величина очередей и достижимость границы определяются полным фазовым профилем. Вычислительный пример такого расхождения приведён в разделе~\ref{sec:computational-results}.
\end{remark}

\begin{definition}[Resource-Augmented Phase Graph]\label{def:resource-graph}
Зафиксируем назначение каждой задачи одному из $P$ агентных потоков, порядок задач внутри каждого потока и порядок обслуживания всех человеческих фаз. Если эти порядки совместимы с исходными зависимостями, построим ориентированный граф $G_R(\omega)$, вершинами которого служат фазы, а весом вершины --- её эффективная длительность с множителем $C(P)$ для агентной или $\gamma(P)$ для человеческой фазы. В граф добавляются дуги четырёх типов:
\begin{enumerate}
  \item между последовательными фазами одной задачи;
  \item от последней фазы предшественника к первой фазе зависимой задачи;
  \item между последовательными человеческими фазами в выбранном порядке обслуживания;
  \item от последней фазы одной задачи к первой фазе следующей задачи в том же агентном потоке.
\end{enumerate}
Набор решений о назначении и ресурсных порядках обозначим $\omega$. Допустимыми считаются только такие $\omega$, для которых $G_R(\omega)$ ацикличен. Через $\Lambda(\omega)$ обозначим длину максимального взвешенного пути в этом графе.
\end{definition}

\begin{proposition}[Exact Representation of the M4 Optimum]\label{prop:resource-graph}
Для фиксированного допустимого ресурсного порядка $\omega$ расписание ранних начал имеет длительность $\Lambda(\omega)$ и является допустимым. Обратно, любое допустимое расписание $\pi$ индуцирует некоторый порядок $\omega(\pi)$, причём $\Lambda(\omega(\pi))\le T(\pi)$. Следовательно,
\[
  T^{*}=\min_{\omega\in\Omega}\Lambda(\omega),
\]
где $\Omega$ --- конечное множество допустимых назначений и ресурсных порядков.
\end{proposition}

\begin{proof}
При фиксированном ациклическом графе время раннего начала каждой фазы равно максимальной длине пути до неё. Дуги первых двух типов обеспечивают технологический порядок, третьего типа --- непересечение человеческих фаз, четвёртого --- непересечение задач, закреплённых за одним потоком. Поэтому построенное расписание допустимо и завершается за $\Lambda(\omega)$. В обратную сторону непересекающиеся операции любого расписания можно упорядочить по каждому ресурсу; все добавленные дуги согласованы с фактическим временем и потому не образуют цикла, а длина любого пути не превосходит $T(\pi)$. Минимизация двух представлений даёт равенство.
\end{proof}

Это представление играет роль \emph{resource-augmented critical path}: исходный критический путь дополняется дугами конкуренции за человека и агентные потоки. Оно не вводит новый уровень M5, поскольку не добавляет ограничение к M4, а раскрывает точную комбинаторную структуру уже заданного расписания. Такая конструкция близка к дизъюнктивным графам и ресурсно-ограниченным моделям расписаний \cite{brucker1999resource}.

\begin{corollary}[Attainability Certificate for the Lower Bound]\label{cor:attainability}
Для любого экземпляра M4 выполняется $B_4\le T^{*}$. Равенство $T^{*}=B_4$ имеет место тогда и только тогда, когда существует допустимое расписание $\pi$ с $T(\pi)=B_4$; эквивалентно, существует $\omega\in\Omega$ с $\Lambda(\omega)=B_4$. Явное расписание такой длины служит сертификатом достижимости. Если решатель доказывает отсутствие расписания с горизонтом $B_4$, граница недостижима и $T^{*}>B_4$.
\end{corollary}

\begin{corollary}[Simple Tightness Cases]\label{cor:simple-tightness}
Если $P=1$ либо task-DAG является полной последовательной цепочкой всех задач,
то $T^*=B_4=W_4$.
\end{corollary}

\begin{proof}
При $P=1$ имеем $C(1)=\gamma(1)=1$, а все задачи последовательно занимают
единственный поток. Их допустимо выполнить в любом топологическом порядке без
простоя, поэтому срок равен сумме эффективных фаз $W_4$; остальные ветви
$B_4$ её не превышают. В полной цепочке никакие две задачи не могут
пересекаться при любом $P$, её длина равна сумме весов $L_4=W_4$, и
последовательное расписание достигает этой границы.
\end{proof}

Достижимость горизонта $B_4$ можно проверять как задачу существования
допустимого расписания, а $T^{*}$ --- прямой минимизацией makespan. Конкретная
вычислительная формулировка и протокол проверки приведены в
разделе~\ref{sec:computational-method}.

\begin{proposition}[Monotonicity of Available Capacity]\label{prop:capacity-monotonicity}
Если длительности фаз и параметры $C$, $\gamma$ зафиксированы и не меняются при добавлении потока, то
\[
  T^{*}(P+1)\le T^{*}(P).
\]
\end{proposition}

\begin{proof}
Любое расписание для $P$ потоков допустимо при $P+1$: дополнительный поток можно не использовать. Поэтому множество допустимых расписаний только расширяется.
\end{proof}

\begin{proposition}[Existence of a Finite Optimal Configuration]\label{prop:finite-optimum}
Пусть сравниваются конфигурации с
\[
 C(P)=1+\alpha(P-1)+\beta P(P-1),\qquad
 \gamma(P)=1+\delta(P-1),
\]
где $P\in\mathbb N$. Если $A>0$ и $\beta>0$ либо $H>0$ и $\delta>0$, то $T^{*}(P)\to\infty$ при $P\to\infty$, а значит минимум $T^{*}(P)$ достигается при конечном $P$.
\end{proposition}

\begin{proof}
В первом случае $W_4(P)/P\ge C(P)A/P\sim\beta AP$; во втором $B_4\ge\gamma(P)H\sim\delta HP$. В обоих случаях $T^{*}(P)\ge B_4(P)\to\infty$. Последовательность на натуральных $P$, уходящая в бесконечность, достигает глобального минимума на конечном начальном отрезке.
\end{proof}

Предложения~\ref{prop:capacity-monotonicity} и~\ref{prop:finite-optimum} относятся к разным вопросам. Первое запрещает ухудшение от неиспользуемой мощности при неизменных длительностях. Второе сравнивает разные организационные конфигурации, в которых рост настроенного параллелизма сам изменяет издержки $C(P)$ и $\gamma(P)$. Поэтому внутренний минимум по $P$ не следует из одной лишь дискретности расписания или очередей, но гарантируется указанными растущими штрафами.

\begin{proposition}[Human-Resource Speedup Ceiling]\label{prop:amdahl}
Для фиксированного набора задач и назначения режимов $\rho$, если $H>0$, при любом параллелизме $P \ge 1$ и любых $C(P) \ge 1$, $\gamma(P) \ge 1$:
\[
  S_E^{\mathrm{ach}}:=\frac{T_h}{T_{ai}}
  \;\le\; \frac{T_h}{\gamma(P)\,H} \;\le\; \frac{T_h}{H} \;=\; \frac{1}{\bar{h}_w},
  \qquad
  \bar{h}_w := \frac{\sum_{i=1}^{N} Z_i h_i^{(\rho(i))}}{\sum_{i=1}^{N} Z_i}.
\]
\end{proposition}

\begin{proof}
Человеческие интервалы всех задач обслуживаются одним разработчиком, поэтому попарно не пересекаются во времени, а их суммарная длительность составляет не менее $H$ (и не менее $\gamma(P) H$ с учётом переключений). Следовательно, $T_{ai} \ge \gamma(P) H \ge H$, откуда $S_E^{\mathrm{ach}} = T_h / T_{ai} \le T_h / H = 1/\bar{h}_w$. Ч.Т.Д.
\end{proof}

При $H=0$ обязательные человеческие фазы отсутствуют, поэтому ограничение
является тривиальным и потолок по человеческому ресурсу условно принимается
равным бесконечности. Остальные ветви $B_4$ и ограничения расписания сохраняются.

\begin{remark}[Analogy with Amdahl's Law]\label{rem:amdahl-analogy}
Предложение~\ref{prop:amdahl} представляет собой структурный аналог закона Амдала: роль последовательной составляющей играет средневзвешенная нормированная человеческая работа $\bar{h}_w$ \cite{amdahl1967validity}. Потолок $1/\bar{h}_w$ не зависит от $P$. Если целевое ускорение превышает эту величину, достичь его добавлением агентов невозможно; необходимо уменьшать человеческую составляющую $h$ за счёт более автономных режимов, автоматизации проверки или формализованных критериев приёмки. На уровне M4 верхняя оценка ускорения по суммарной занятости потоков равна $P/[C(P)\bar a_w+\gamma(P)\bar h_w]$, а по человеческому ресурсу --- $1/[\gamma(P)\bar h_w]$. Эти два агрегата полезны для быстрой диагностики, но не учитывают критический путь и очередь к человеку.
\end{remark}

\subsection{Cross-Level Properties}

При согласованной параметризации последовательные уточнения сохраняют
ограничения предыдущих уровней. Поскольку $M_N\le L$, $C(P)\ge1$ и
$\gamma(P)\ge1$, имеем $W_3(P)\ge W$, $L_3(P)\ge L$,
$W_4(P)\ge W_3(P)$ и $L_4(P)\ge L_3(P)$. Поэтому нижние оценки упорядочены
следующим образом:
\[
  B_0 \le B_1 \le B_2 \le B_3 \le B_4,
\]
Поэтому верхние оценки достижимого ускорения $\bar S^{(j)} := T_h / B_j$ монотонно убывают с ростом $j$. На уровне M0 время в точности совпадает с оценкой ($T_{ai}^{(0)} = B_0$), а Наблюдение~\ref{obs:ideal} задаёт алгебраически эквивалентный критерий ускорения в форме необходимого и достаточного условия. Начиная с уровня M1 величина $B_j$ служит только нижней границей. Фактическое время определяется распределением задач между исполнителями и может строго превышать $B_j$ даже при оптимальном расписании; пример такого расхождения для $B_1$ приведён в разделе с расчётами. Условие ускорения по линейной составляющей при этом остаётся необходимым, но уже не является достаточным (Замечание~\ref{rem:threshold}).

Связь с обозначениями вероятностного расширения модели: линейной части уровня M3 соответствует $T_{ai}^{\mathrm{lin}}:=W_3(P)/P$, а коррекция на неделимость требует использовать максимальный вес $\max_i w_i^{(3)}(P)$. Граф зависимостей, фазовая очередь и общий человеческий ресурс в вероятностном расширении пока не рассматриваются.

\subsection{Scale Invariance and Comparative Statics}

\begin{proposition}[Scale Invariance of Configuration Comparison]\label{prop:invariance}
Зафиксируем декомпозицию проекта --- набор задач со сложностями $Z_i$ и граф
зависимостей $G$ --- и конфигурацию $\sigma$ с целым числом исполнителей $P$.
На уровне M4 дополнительно фиксируются для каждой задачи число, типы и порядок
фаз, отношения предшествования, правило локальной блокировки, $C(P)$ и
$\gamma(P)$. Масштабное преобразование умножает длительность \emph{каждой}
агентной и человеческой фазы на один общий коэффициент $\kappa>0$; следовательно,
$a_i$ и $h_i$ масштабируются по отдельности, а их разбиение и фазовый профиль
сохраняются. Пусть $T^{*}(\kappa)$ --- оптимальное время завершения после такого
преобразования. Тогда:
\begin{enumerate}
  \item $T^{*}(\kappa) = \kappa\, T^{*}(1)$ (однородность первой степени); тем же свойством обладают нижние оценки $B_0$--$B_4$, если $C(P)$ и $\gamma(P)$ от $\kappa$ не зависят. Оптимальное распределение задач между исполнителями, порядок обслуживания очереди, активный член $\max\{W_4(P)/P,\,L_4(P),\,\gamma(P)H\}$ и состав критического пути от $\kappa$ не зависят;
  \item если для двух вариантов организации работы $A$ и $B$ полные векторы длительностей соответствующих агентных и человеческих фаз известны с точностью до одного и того же общего множителя $\kappa>0$, то есть $\mathbf d^{A}=\kappa\,\hat{\mathbf d}^{A}$ и $\mathbf d^{B}=\kappa\,\hat{\mathbf d}^{B}$, то отношение $T^{*}_A/T^{*}_B$ и, следовательно, ранжирование вариантов не зависят от $\kappa$. В частности, $a_i$, $h_i$ и $k_i=a_i+h_i$ должны масштабироваться согласованно; одной пропорциональности суммарных $k_i$ при меняющемся разбиении на фазы недостаточно.
\end{enumerate}
\end{proposition}

\begin{proof}
При умножении длительностей всех фаз на $\kappa$ длительности задач, человеческое обслуживание и возникающие при том же порядке событий интервалы ожидания умножаются на $\kappa$. Любое допустимое расписание преобразуется в допустимое расписание новой задачи умножением всех моментов начала и завершения на $\kappa$. Порядок фаз, отношения предшествования, локальное закрепление агентов и ограничения ресурсов сохраняются. Преобразование обратимо и задаёт взаимно однозначное соответствие между множествами допустимых расписаний. Поэтому $T^{*}(\kappa)=\kappa T^{*}(1)$, а оптимальное расписание переходит в оптимальное. Однородность $B_j$ следует из линейности $W$, $M_N$, $L$, $W_3$, $L_3$, $W_4$, $L_4$ и $H$ по $\kappa$. Утверждение (ii) следует из (i): $T^{*}_A(\kappa)/T^{*}_B(\kappa)=T^{*}_A(1)/T^{*}_B(1)$. Ч.Т.Д.
\end{proof}

\begin{remark}[Comparative Statics without Exact Calibration]\label{rem:comparative}
Предложение~\ref{prop:invariance} уточняет возможности модели. Абсолютный прогноз --- значение $T_{ai}$ в часах или $S_E$ --- требует знания уровня длительностей и в том же масштабе наследует ошибку калибровки (Замечание~\ref{rem:estimation}). Однако \emph{выбор между вариантами декомпозиции и организации работы} не зависит от общего масштаба, если полные относительные профили агентных и человеческих фаз известны с точностью до одного общего множителя.

Направление выгодных изменений модель также указывает без точной калибровки. Нижняя оценка $B_2$ кусочно-линейна по вектору $k = (k_1, \dots, k_N)$. Если активна ветвь $L>W/P$ и критический путь $\pi^*$ единственен, её эластичности по отдельным коэффициентам вычисляются явно:
\[
  \frac{\partial B_2}{\partial k_i} \cdot \frac{k_i}{B_2} =
  \begin{cases}
    \dfrac{Z_i k_i}{\sum_{j} Z_j k_j}, & \text{если } W/P > L, \\[2.5ex]
    \dfrac{Z_i k_i X}{L}, & \text{если } L > W/P,\ \pi^*\text{ единственен и }i \in \pi^{*}, \\[2.5ex]
    0, & \text{если } L > W/P,\ \pi^*\text{ единственен и }i \notin \pi^{*},
  \end{cases}
\]
В точках переключения активных членов и при нескольких критических путях $B_2$ может быть недифференцируема; тогда используются направленные производные, определяемые всеми активными путями. При активной ветви $L$ увеличение $P$ не снижает саму нижнюю оценку $B_2$, однако срок $T^*$ может уменьшиться, пока не достигнет этой границы. Аналогично, при активной ветви $\gamma(P)H$ доля задачи в $H$ описывает чувствительность $B_4$, но переносить её на $T^*$ можно только при доказанной тесноте $T^*=B_4$.

Инвариантность имеет три ограничения. Во-первых, она относится только к
\emph{общей мультипликативной} ошибке всех фаз. Если ошибки оценивания искажают
относительные длительности задач, разбиение между $a_i$ и $h_i$ либо внутренний
фазовый профиль, состав критического и ресурсно-дополненного путей и ранжирование
вариантов могут измениться. Во-вторых, множитель $\kappa$ должен быть общим для
всех фаз и всех сравниваемых вариантов. Иными словами, систематическое смещение
процедуры оценки (Замечание~\ref{rem:estimation}) предполагается одинаковым для
конфигураций одной команды при едином протоколе измерения. В-третьих, сравнение
разных режимов требует знания относительных режимных профилей отдельно для
агентных и человеческих фаз. Это требование существенно слабее знания их
абсолютных длительностей, но сильнее знания только суммарных коэффициентов $k_i$.
\end{remark}

\subsection{Calibration and Uncertainty}

\begin{remark}[Stochastic Nature of the Coefficients $k_i$]\label{rem:stochastic}
На практике влияние ИИ на выполнение задачи недетерминировано. Коэффициент $k_i$ целесообразно рассматривать как случайную величину, распределение которой отражает неопределённость генерации кода, вероятность дефектов и затраты времени на их исправление. Тогда $T_{ai}$ также становится случайной величиной, а условие ускорения следует проверять для математического ожидания $\mathbb{E}[T_{ai}] < T_h$ или выбранных квантилей распределения.

Отметим, что в силу линейности математического ожидания $\mathbb{E}[T_{ai}^{(0)}] = \frac{X}{P}\sum_{i=1}^{N} Z_i\,\nu_i$, где $\nu_i = \mathbb{E}[k_i]$, поэтому тождественный критерий уровня M0 для детерминированных коэффициентов (Наблюдение~\ref{obs:ideal}) переносится на ожидания дословно:
\[
  \mathbb{E}[T_{ai}^{(0)}] < T_h \;\iff\; P > \frac{\sum_{i=1}^{N} Z_i \nu_i}{\sum_{i=1}^{N} Z_i}.
\]
Систематическое развитие стохастической постановки (дисперсия, квантильные гарантии) вынесено в отдельную вероятностную модель.
\end{remark}

\begin{remark}[Estimating $k_i$ before Project Start]\label{rem:estimation}
Операциональное определение $k_i^{(r)}$ (Определение~\ref{def:params}) носит ретроспективный характер: коэффициент вычисляется по уже наблюдавшемуся времени $T_i^{ai}(r)$ и сам по себе не является прогнозом. Чтобы применять критерии ускорения (Наблюдение~\ref{obs:ideal}, Замечание~\ref{rem:threshold}, Теорема~\ref{thm:sufficient}) \emph{до} начала проекта, требуется внешняя оценка $\hat{k}_i^{(r)}$. Её можно получить следующим образом:
\begin{enumerate}
  \item задачи классифицируются по типам (CRUD-операции, интеграция с внешним API, тесты, инфраструктура и т.\,п.); тип задачи $i$ обозначим $c(i)$;
  \item по задачам команды накапливается выборка значений $k$ для пар «тип $\times$ режим»; раздельный хронометраж позволяет одновременно фиксировать человеческую составляющую $h$ (Определение~\ref{def:params}). Неуспешные запуски, тайм-ауты и непринятые результаты должны учитываться явно либо как наблюдения, либо как цензурированные случаи: выборка только завершённых задач оценивает условную длительность и подвержена смещению выжившего;
  \item оценка $\hat{k}_i^{(r)}$ берётся как эмпирическое среднее по паре $(c(i), r)$ --- для критерия в терминах ожиданий (Замечание~\ref{rem:stochastic}) --- либо как верхний квантиль выборки для консервативного решения.
\end{enumerate}
На уровне M0 чувствительность результата к ошибке оценки следует непосредственно из $S_E = P/\bar{k}_w$: относительная ошибка $\bar{k}_w$ в том же масштабе переносится в $S_E$ (завышение $\bar{k}_w$ в $(1+\varepsilon)$ раз занижает $S_E$ в $(1+\varepsilon)$ раз). Для сравнения конфигураций на M0 абсолютный уровень $k$ менее важен: общая мультипликативная ошибка не меняет их ранжирование. На M4 тот же вывод требует более сильной предпосылки Предложения~\ref{prop:invariance}: один множитель должен масштабировать все агентные и человеческие фазы, сохраняя их типы и порядок. Следовательно, модель имеет двойной статус. При ретроспективном применении она задаёт точное тождество, а при использовании для принятия решений точность абсолютных прогнозов $T_{ai}$ и $S_E^{\mathrm{ach}}$ определяется качеством оценки $\hat{k}$ и фазового профиля. Сравнительные выводы устойчивы только к общей систематической ошибке в смысле Замечания~\ref{rem:comparative}.
\end{remark}

\ifdefined\FULLARTICLE
  \let\finishdocument\relax
\else
  \def\finishdocument{\end{document}}
\fi
\finishdocument
